% TOP_PROBLEM: {"id": 30006790, "title": "Effective Roth Theorem", "category_id": 1, "category": {"id": 1, "name": "number_theory", "display_name": "Number Theory", "description": "Properties of integers, prime numbers, Diophantine equations.", "slug": "number-theory", "order_index": 1, "created_at": "2026-07-31T15:26:25.670Z"}, "status": "open"}
% FIRST_POSTED: 2026-09-27
% Standalone research entry 208; batch 208--217.
% Original section material preserved from Pasted text(3).txt.
% Research additions: 27 September 2026.
% Compile twice with: lualatex 30006790.tex
% !TeX program = lualatex
% Compile twice: lualatex open_problems_500.tex
% All references and prose are embedded; no BibTeX database or external inputs.
\documentclass[11pt,a4paper]{article}
\usepackage[margin=24mm,headheight=14pt]{geometry}
\usepackage{fontspec}
\setmainfont{Latin Modern Roman}
\setsansfont{Latin Modern Sans}
\setmonofont{Latin Modern Mono}
\usepackage{amsmath,amssymb,mathtools}
\usepackage{unicode-math}
\setmathfont{Latin Modern Math}
\usepackage{microtype}
\usepackage{enumitem,xurl,needspace}
\usepackage[dvipsnames]{xcolor}
\usepackage{hyperref}
\hypersetup{unicode=true,colorlinks=true,linkcolor=MidnightBlue,urlcolor=MidnightBlue,
 pdftitle={500 Mathematical Problem Statements},pdfauthor={Source-annotated compilation},bookmarksnumbered=true}
\usepackage{bookmark}
\usepackage{tocloft}
\setlength{\cftsecnumwidth}{3em}
\cftsetpnumwidth{2.5em}
\cftsetrmarg{4em}
\usepackage{titlesec}
\titleformat{\section}{\Large\bfseries\raggedright}{\thesection}{0.7em}{}
\usepackage{fancyhdr}
\pagestyle{fancy}\fancyhf{}
\fancyhead[L]{\small\sffamily Mathematical problem statements}
\fancyhead[R]{\small Source edition 22}
\fancyfoot[C]{\thepage}
\setlength{\parindent}{0pt}
\setlength{\parskip}{5pt plus 1pt}
\setlength{\emergencystretch}{3em}
\setcounter{tocdepth}{1}
\setcounter{secnumdepth}{1}
\setlist[itemize]{leftmargin=3em,itemsep=5pt,topsep=4pt,parsep=0pt}
\allowdisplaybreaks
\newcommand{\N}{\mathbb{N}}
\newcommand{\Z}{\mathbb{Z}}
\newcommand{\Q}{\mathbb{Q}}
\newcommand{\R}{\mathbb{R}}
\newcommand{\C}{\mathbb{C}}
\newcommand{\F}{\mathbb{F}}
\newcommand{\A}{\mathbb{A}}
\newcommand{\PP}{\mathbb P}
\newcommand{\E}{\mathbb{E}}
\newcommand{\eps}{\varepsilon}
\newcommand{\dd}{\,\mathrm d}
\newcommand{\sref}[2]{\hyperlink{source:#1:#2}{[S#2]}}
\newcommand{\eref}[2]{\hyperlink{extra:#1:#2}{[E#2]}}
% Turn the next switch off for a shorter reading copy. The quotations preserve
% every exactTarget field, including qualifications omitted from a short summary.
\newif\ifshowcatalogue
\showcataloguetrue
\newcommand{\cataloguescope}[1]{\ifshowcatalogue
 \paragraph{Exact scope in the supplied catalogue.}
 \begin{quote}\small #1\end{quote}\fi}

% Standalone wrapper; no external inputs or bibliography files are required.
\fancyhead[L]{\small\sffamily Research notebook: section 30006790}
\fancyhead[R]{\small 27 September 2026}
\hypersetup{pdftitle={Research attempt 208},pdfauthor={Research notebook}}
\setlength{\parskip}{4pt plus 1pt}
\setlist[itemize]{before=\small\interlinepenalty10000,itemsep=3pt}
\begin{document}
\setcounter{section}{0}
% ================================================================
% problemId: 30006790
% Canonical problem ID: 30006790
% exactTarget SHA-256: 0add024928cac15703b995dbcb3b162a968b879b0cc4c6d5b1b4ba440426a52d
\clearpage
\section*{\texorpdfstring{Effective Roth Theorem}{Effective Roth Theorem}}\label{30006790}
\subsection{Definitions and mathematical statement}
Encode a real algebraic irrational $\alpha$ by its integer minimal polynomial and a rational interval containing exactly the chosen real root. Given also rational $\eps>0$, the requested single total algorithm must output an integer $Q\ge1$ such that
\[
 \gcd(p,q)=1,\quad q\ge1,\quad
 |\alpha-p/q|<q^{-2-\eps}\quad\Longrightarrow\quad q\le Q.
\]
The algorithm must halt on every valid input; $p$ may be any integer. The bound can be enormous and the running time unrestricted. Merely proving that finitely many such approximations exist, or giving a bound depending on an unspecified noncomputable constant, does not supply this effective version.
\par\smallskip\textit{Formulation and concept references:} \sref{30006790}{1}.
\cataloguescope{There exists one total computable algorithm which, given a real algebraic irrational alpha by its integer minimal polynomial and a rational isolating interval, and a positive rational epsilon, returns an integer Q>=1 such that every reduced rational p/q, with p an integer and q>=1 an integer, satisfying |alpha-p/q|<q\textasciicircum{}(-2-epsilon) has q<=Q. The algorithm must halt with this guarantee on every valid input. No time, space, practical-size or software-implementation requirement is imposed.}
\subsection{Short English statement}
Can an algorithm compute a guaranteed denominator cutoff for every exceptionally accurate rational approximation to a specified algebraic irrational?
% BEGIN RESEARCH ADDITION 208
\subsection{Research attempt}

\paragraph{Current frontier.}
The issue is a computable location of the last exceptional denominator, not
finiteness of the exceptional set. Dill's introduction distinguishes the
effective Liouville exponent and special effective improvements from Roth's
ineffective exponent $2+\eps$.\sref{30006790}{1} Adiceam--Shirandami give effective
\emph{proportions} of algebraic inputs with approximations in specified
ranges; their order of limits and height regime do not give a cutoff for
each fixed input.\eref{30006790}{1} The argument below does not use such a density
statement as a pointwise theorem.

\paragraph{Attack route.}
There are three natural places to seek a stopping certificate: a polynomial
lower bound, the continued-fraction tail, and exhaustion of a quantitatively
finite exceptional set. The first works beyond the Liouville threshold.
The second isolates a specific infinite-tail estimate. The third is tested
against a computability construction in which a single late exceptional
convergent encodes a halting time. This last test deliberately changes the
\emph{representation} of the input; it identifies information a successful
proof must use, rather than disproving the stated problem.

\paragraph{Attempt.}
\textbf{1. An elementary, fully effective region.}
Write the supplied irreducible polynomial as
$P(X)=\sum_{j=0}^{d}a_jX^j$, with $d\ge2$.
From the isolating interval choose an integer $R>|\alpha|+1$ and compute
\[
 L=\sum_{j=1}^{d}j|a_j|R^{j-1}>0.
\]
For a candidate satisfying the target inequality, $|\alpha-p/q|<1$,
so both $\alpha$ and $p/q$ lie in $(-R,R)$. Irreducibility implies
$P(p/q)\ne0$, while $q^dP(p/q)$ is an integer. The mean-value theorem gives
\begin{equation}\label{30006790:liouville}
 q^{-d}\le |P(p/q)|\le L|\alpha-p/q|,
 \qquad |\alpha-p/q|\ge L^{-1}q^{-d}.
\end{equation}
Consequently, if $\delta=2+\eps-d>0$, an exception must satisfy
$q^\delta<L$. An integer $Q\ge1$ with $Q^\delta\ge L$ can be found by
integer search and exact comparison of rational powers. It is a valid
cutoff. In particular, this supplies a total algorithm for every quadratic
input and every positive rational $\eps$; no continued-fraction period has
to be computed. For $d\ge3$ and $0<\eps\le d-2$, the displayed inequality
has the wrong exponent and gives no bound on $q$.

\textbf{2. A precise tail certificate.}
Let $p_n/q_n$ be the regular continued-fraction convergents of $\alpha$.
For sufficiently large $n$, their positive denominators strictly increase.
The determinant relation for adjacent convergents and the complete
quotient formula give
\begin{equation}\label{30006790:cf}
 \frac{1}{q_n(q_{n+1}+q_n)}
 <\left|\alpha-\frac{p_n}{q_n}\right|
 <\frac{1}{q_nq_{n+1}}.
\end{equation}
Indeed, writing the next complete quotient as $a_{n+1}+\theta$, with
$0<\theta<1$, the exact denominator of the error is
$q_n(q_{n+1}+\theta q_n)$.
If one could compute $n_0$ for which
\begin{equation}\label{30006790:tail}
 q_{n+1}+q_n\le q_n^{1+\eps}\qquad(n\ge n_0),
\end{equation}
then no such convergent would be exceptional. A reduced rational with
$q>2^{1/\eps}$ and error $<q^{-2-\eps}$ has error $<1/(2q^2)$ and hence
is a convergent, by the elementary continued-fraction best-approximation
criterion. Thus a bound covering $2^{1/\eps}$ and the finitely many earlier
convergent denominators would solve the input instance.

The problem is not that the tail inequality is false. Apply Roth with
$\eps/2$ and use the upper bound in \eqref{30006790:cf}: eventually
$q_{n+1}<q_n^{1+\eps/2}$, which implies \eqref{30006790:tail} once
$q_n^{\eps/2}$ is sufficiently large. This argument proves existence of a
tail but supplies no computable starting index. Importing this index from
Roth would be exactly the effectivity gap.\sref{30006790}{1}

\textbf{3. Exact enumeration and the gap principle do not certify a tail.}
For each $q$, only finitely many integers $p$ with $|p|<Rq$ need examination.
The comparison of the algebraic number $|\alpha-p/q|$ with the positive
algebraic number $q^{-2-\eps}$ is decidable, including equality. Enumeration
therefore gives a computable nondecreasing sequence of provisional maxima.
Roth implies that this sequence eventually stabilizes, but stabilization
is not an observable event in a finite prefix.

For two distinct exceptional reduced fractions with $q'\ge q$,
\[
 \frac1{qq'}\le\left|\frac pq-\frac{p'}{q'}\right|
 <q^{-2-\eps}+(q')^{-2-\eps}\le2q^{-2-\eps},
 \qquad q'>\tfrac12q^{1+\eps}.
\]
This is a \emph{lower} separation bound. It permits an arbitrarily late
next exception; it is not an upper bound on the last one. Even a computable
upper bound on the \emph{number} of exceptions does not say when enumeration
has finished if the actual number is smaller.

\textbf{4. A representation-sensitive impossibility lemma.}
Consider a different input format: a program producing arbitrarily accurate
rational approximations to a real number, promised to be quadratic
irrational. There is no total computable cutoff procedure on all such
programs even for the fixed exponent $\eps=1$.

Here is a direct proof. Given a machine index $e$, construct a continued
fraction $\alpha_e=[0;a_1,a_2,\ldots]$ as follows. All digits are $1$ unless
the machine first halts at step $t\ge0$. In that event put $n=t+3$, change
just the digit $a_{n+1}$ to $q_n^2+1$, and leave every other digit equal to
$1$. The denominator $q_n$ before this exceptional digit is the denominator
of $[0;1,\ldots,1]$, hence a known increasing Fibonacci number with
$q_n\ge3$.

Every finite digit prefix is computable by a finite simulation. Moreover,
continued-fraction cylinders of length $j$ have diameter at most the
reciprocal square of their least possible denominator, which grows at least
as a Fibonacci number. This gives a \emph{uniform computable Cauchy name}
for $\alpha_e$, with no decision about eventual halting. Every $\alpha_e$
is quadratic irrational: its tail is all ones, whose complete quotient
$\tau$ satisfies $\tau^2-\tau-1=0$, and a finite continued-fraction prefix
acts on $\tau$ by a rational fractional-linear transformation of determinant
$\pm1$.

If the machine halts, the exceptional convergent satisfies
\[
 q_{n+1}=(q_n^2+1)q_n+q_{n-1}>q_n^3,
 \qquad
 \left|\alpha_e-\frac{p_n}{q_n}\right|<q_n^{-4}<q_n^{-3}.
\]
Suppose a cutoff algorithm on the Cauchy-program representation returned
$Q_e$. Choose a computable $B$ such that the all-ones denominator
$q_{B+3}>Q_e$, and simulate the machine for $B$ steps. If it has not halted,
it can never halt later: halting at $t>B$ would create the above exception
with denominator exceeding $Q_e$. This decides the halting problem, a
contradiction.

The elementary undecidability used here follows by diagonalization: a
hypothetical procedure deciding whether a program halts on its own index
would allow a program which halts precisely when that procedure predicts
that it does not. No unproved Diophantine conjecture enters this lemma.

\paragraph{Stress test.}
The representation lemma is \emph{not} a counterexample to Effective Roth.
The original input includes a minimal polynomial and a root-isolating
interval, whereas the constructed family supplies only a Cauchy program.
There is no uniform method to extract that polynomial and interval from
these programs: composing one with the quadratic algorithm in part~1
would give the impossible procedure in part~4. Once a particular machine
has halted, its finite exceptional prefix does determine such a polynomial;
knowing that this has happened is exactly the unavailable information.

No asymptotic assertion was used to terminate a search. The rational-power
comparisons in part~1 and finite algebraic comparisons in part~3 are exact.
The strict inequality in the problem is respected at the cutoff and at
possible equality cases. The continued-fraction argument works after
subtracting an integer as well, so negative $\alpha$ and arbitrary integer
numerators are not excluded. The algorithmic impossibility argument is an
obstruction to discarding the input's algebraic certificate, not evidence
that the correctly encoded target is false.

\paragraph{Outcome.}
\textbf{Outcome: Exploratory approach with unresolved gap.}

The notebook proves the effective region $\eps>d-2$, identifies the missing
computable continued-fraction tail, and proves a representation barrier
against replacing the minimal-polynomial input by unrestricted numerical
access. These are not claimed as new frontier results. For degree at least
three in the remaining exponent range, the missing step is a way to use
the supplied algebraic certificate to bound the last large partial quotient
without importing an ineffective Roth constant. No such step is established.
% END RESEARCH ADDITION 208
\subsection{Sources}
\begin{itemize}\raggedright
\item[\textnormal{[S1]}]\hypertarget{source:30006790:1}{} Effective Approximation and Diophantine Applications; arXiv1601.02243v1,10January2016.
\url{https://arxiv.org/pdf/1601.02243v1}.
{\small\textit{Catalogue formulation source.}}
% BEGIN RESEARCH REFERENCES 208
\item[\textnormal{[E1]}]\hypertarget{extra:30006790:1}{} Faustin Adiceam and Victor Shirandami,
\emph{Probabilistic Effectivity in the Subspace Theorem},
arXiv:2411.01247v2 (13 November 2024), Sections 1.1--1.2, Theorem 1.3
and the order of limits in (1.9). Research in Number Theory 12, article 8 (2026),
DOI 10.1007/s40993-025-00692-0.
\url{https://arxiv.org/abs/2411.01247}.
{\small\textit{Consulted 27 September 2026. Used for the probabilistic/pointwise distinction, not as a pointwise effective Roth theorem.}}
% END RESEARCH REFERENCES 208
\end{itemize}

\end{document}
